% ============================================================
% ARQUIVO: src/config.tex
% DESCRIÇÃO: Configurações globais do template LaTeX
% ATENÇÃO: Este arquivo NÃO deve ser editado pelo usuário!
% ============================================================

% --- PACOTES ESSENCIAIS ---
\usepackage[utf8]{inputenc}
\usepackage[T1]{fontenc}
\usepackage[brazil]{babel}
\usepackage{geometry}
\usepackage{setspace}
\usepackage{amsmath, amssymb}
\usepackage{graphicx}
\usepackage{indentfirst}
\usepackage{microtype}
\usepackage{xcolor}
\usepackage{ifthen}
\usepackage{hyperref}
\usepackage{float}
\usepackage{caption}
\usepackage{booktabs}
\usepackage{array}
\usepackage{longtable}
\usepackage{enumitem}

% --- CONFIGURAÇÕES DE MARGEM ---
\geometry{
    a4paper,
    top=30mm,
    bottom=20mm,
    left=30mm,
    right=20mm
}

% --- ESPAÇAMENTO ---
\onehalfspacing

% --- CORES INSTITUCIONAIS ---
\definecolor{verdeUnesp}{RGB}{0, 102, 51}
\definecolor{azulUnesp}{RGB}{0, 51, 102}
\definecolor{cinzaUnesp}{RGB}{128, 128, 128}

% --- HYPERREF ---
\hypersetup{
    colorlinks=true,
    linkcolor=azulUnesp,
    urlcolor=azulUnesp,
    citecolor=verdeUnesp,
    bookmarks=true,
    pdfpagemode=UseNone
}

% --- LEGENDAS ---
\captionsetup{
    font=small,
    labelsep=endash,
    justification=centering
}

% --- LISTAS ---
\setlist{
    topsep=6pt,
    itemsep=2pt,
    partopsep=0pt,
    parsep=0pt,
    leftmargin=2.5em
}

% --- TABELAS ---
\newcolumntype{L}[1]{>{\raggedright\arraybackslash}p{#1}}
\newcolumntype{C}[1]{>{\centering\arraybackslash}p{#1}}
\newcolumntype{R}[1]{>{\raggedleft\arraybackslash}p{#1}}

% ============================================================
% FIM DO ARQUIVO config.tex
% ============================================================
